\thispagestyle{plain}
% \chapter*{Resumo}

% \guiainfo{Na secção \textit{Resumo}, apresente um resumo conciso do seu projeto, destacando os pontos principais. Comece com uma breve declaração do problema ou objetivo, seguido de uma descrição da sua abordagem ou metodologia. Resuma os principais resultados ou conclusões, salientando a sua importância ou implicações. Conclua com uma ou duas frases sobre a contribuição global ou o impacto do seu trabalho. O resumo deve ser claro e conciso, idealmente com 150-250 palavras, para que os leitores compreendam rapidamente o seu trabalho e a sua importância.}

% \keywordspt{Palavra-Chave A, Palavra-Chave B, Palavra-Chave C.}

%\MediaOptionLogicBlank

\pdfbookmark[1]{Abstract}{abstract}
\chapter*{Abstract}
% \guideinfo{In the \textit{Abstract} section, provide a concise summary of your project, highlighting the key points. Begin with a brief statement of the problem or objective, followed by a description of your approach or methodology. Summarise the main results or findings, emphasising their significance or implications. Conclude with a sentence or two on the overall contribution or impact of your work. Keep the abstract clear and focused, ideally within 150-250 words, to give readers a quick understanding of your research and its importance.}

{
Waste management is an urgent environmental and public health challenge across many African urban and peri-urban centers. This report details the design, methodology, and implementation of "Ayang," a web application developed to address this issue. Ayang, named after the Ibibio word for broom, combines a community-focused platform with AI to make trash collection more engaging, encourage civic involvement, and provide useful data for waste management authorities. The system uses Google's Gemini Vision model to analyze user-submitted images, estimating the weight and composition of trash and verifying its collection. By awarding points for successful cleanups, Ayang incentivizes community participation and creates a scalable, data-driven solution to a pressing local problem. This report covers the project's architecture, the AI methodology including plans for fine-tuning, potential performance metrics, and a roadmap for future development and deployment. We argue that by combining gamification with robust AI, Ayang presents a novel and effective model for sustainable, community-led environmental stewardship.}

\keywordsen{Waste Management, Community Engagement, AI, Gamification, Environmental Stewardship.}

%\MediaOptionLogicBlank